\documentclass[11pt]{article}
\usepackage[margin=25mm]{geometry}
\usepackage{fontspec}
\setmainfont{Latin Modern Roman}
\usepackage{amsmath,amssymb,mathtools}
\usepackage{unicode-math}
\setmathfont{Latin Modern Math}
\usepackage{xurl,hyperref}
\hypersetup{colorlinks=true,urlcolor=blue}
\setcounter{secnumdepth}{0}
\setlength{\parindent}{0pt}
\setlength{\parskip}{5pt}
\setlength{\emergencystretch}{3em}
\newcommand{\sref}[2]{\hyperlink{source:#1:#2}{[S#2]}}
\newcommand{\cataloguescope}[1]{\paragraph{Recorded target.}#1}
\newcommand{\R}{\mathbb R}
\newcommand{\C}{\mathbb C}
\newcommand{\N}{\mathbb N}
\newcommand{\Z}{\mathbb Z}
\newcommand{\Q}{\mathbb Q}
\newcommand{\F}{\mathbb F}
\newcommand{\eps}{\varepsilon}
\begin{document}
% UnsolvedMath ID 20002511; source catalogue code AIM-PHYSICS-0019
\setcounter{section}{20002510}
\section{Hurwitz class numbers, CM enumeration, and local conductor layers}\label{20002511}
{\small UnsolvedMath ID 20002511 · Mathematical Physics}
\subsection{Definitions and mathematical statement}

\paragraph{Shared notation.}$\mathbb N=\{1,2,\ldots\}$; the standard number fields have their usual meanings. All additional conventions are stated locally.


\paragraph{Definitions and notation.}A positive-definite integral binary quadratic form is $ax^2+bxy+cy^2$, with discriminant $b^2-4ac<0$, considered up to $SL_2(\mathbb Z)$ equivalence. Hurwitz class numbers count such classes with automorphism weights, including $1/2$ and $1/3$ for the extra-symmetry classes in the source’s convention. The convention $h_0=-1/12$ is separate from the positive-index class-number counting convention.

\paragraph{Catalogue wording retained for traceability.}
With the source’s Hurwitz class-number convention, including $h_0=-1/12$, determine the enumerative significance of the sequence $h_m$ and how to find $h_m$ for $m\ge9$. The source points to the Bringmann--Kane, Bringmann--Murthy and Bringmann--Manschot works for connections between number theory and physics.
\paragraph{Scope and source qualification.}The source says coefficients of h_m even though h_m already denotes coefficients; this is treated as asking about the sequence. Ordinary finite enumeration is not in itself an unresolved obstruction. The intended additional physical enumerative interpretation remains broad; cited papers were not independently reviewed for current status. The catalogue title/status was paired with a machine-generated research attempt; that attempt is not peer-reviewed evidence of settlement. The statement covers the original source questions rather than only the title’s subcase.
\subsection{Short English statement}
Explain the enumerative meaning of Hurwitz class numbers and compute the values requested in the source.
\subsection{Research attempt}
\paragraph{Model, date and investigated scope.}
OpenAI GPT-6 Astra (configured model \texttt{gpt-6-astra}), ultra reasoning; generated 11 October 2026. Computed h9=h10=0, h11=1 and h12=4/3 by reduced-form enumeration, with imprimitive and automorphism weights. The requested broader physical enumerative significance remains unresolved.

\paragraph{Sources and literature check.}
Read the supplied complete definition, catalogue wording, scope qualifications and source list for this canonical ID. The original-source citations are inherited as supplied references; no fresh primary-source or novelty search was performed for this attempt. No current open-status claim is inferred from this calculation.

\paragraph{Finite reduction and the first requested values.}Use the conventional weight \(2/|\operatorname{Aut}_{SL_2}(Q)|\) for a positive-definite integral form \(Q=(a,b,c)\) of discriminant \(-m\). The substitution \(x\mapsto x+ky\) lets one reduce the middle coefficient relative to the leading one after exchanging variables as needed. More explicitly, choose \(k\) to put \(-a<b\le a\); if \(c<a\), exchange the two variables with determinant-one map \((x,y)\mapsto(-y,x)\), which strictly decreases the positive leading coefficient. Iteration terminates with \(|b|\le a\le c\). Then
\[m=4ac-b^2\ge3a^2,\qquad 1\le a\le\sqrt{m/3}.
\]
Thus only finitely many reduced candidates occur for each positive \(m\).
\paragraph{Explicit low-index enumeration.}Discriminants satisfy \(-m\equiv0,1\pmod4\), so \(h_9=h_{10}=0\). For \(m=11\), the bound gives \(a=1\), and the only candidates are \((1,\pm1,3)\), equivalent by an integral translation; there is one class with weight one. For \(m=12\), the candidates are \((1,0,3)\) and \((2,\pm2,2)\), with the latter pair equivalent. Their coefficient gcds differ, so the two classes are inequivalent. The first has only \(\pm I\) as determinant-one automorphisms: its shortest primitive vectors are \(\pm(1,0)\), forcing the remaining basis vector. The second is twice the hexagonal form, whose six shortest vectors give exactly six orientation-preserving symmetries. Hence
\[h_{11}=1,\qquad h_{12}=1+\tfrac13=\tfrac43.
\]
\paragraph{Scope and gap.}This reconstructs elementary arithmetic enumeration, including imprimitive forms, and keeps \(h_0=-1/12\) separate. It does not supply the desired additional physical/geometric enumerative interpretation or audit the cited papers. A concrete moduli problem with these automorphism weights and its counting correspondence is the first missing step. Finite computability is not claimed to be the unresolved mathematical obstacle.

\paragraph{Verification and adversarial audit.}
Checked discriminant congruence, reduction bound, exhaustive small cases, gcd inequivalence, hexagonal weight and separate h0 convention. This was an LLM adversarial audit, not independent human verification; no theorem prover was used. Prompts A then B applied to this ID: Checked discriminant congruence, reduction bound, exhaustive small cases, gcd inequivalence, hexagonal weight and separate h0 convention. Scope remains limited: Identify and prove a geometric or physical moduli-count correspondence with the same Hurwitz weights.

Coordinator checked the reduced-form bound, all candidates at discriminants -11 and -12, their weights and the vanishing inadmissible discriminants. This additional check was model-assisted, not independent human verification.

\paragraph{Tools, process and contributions.}
Python 3 and local shell used to read supplied files and serialize this manuscript; mathematics derived symbolically in the text, with no numerical experiment or formal proof assistant. Contribution prompt A was followed by prompt B. GPT-6 Astra with ultra reasoning drafted the research and performed its stated self-audit; any additional coordinator contribution is identified in the verification and audit. A GPT-6-based Codex coordinator assembled the manuscript and checked repository integration; its exact variant and reasoning configuration were not exposed. No human mathematical contribution, independent human proof check or formal proof certification is supplied. Complete derivations and the consulted source references are included here; private reasoning transcripts are not published.

\paragraph{Final status and first remaining gap.}
\textbf{UNRESOLVED}. Identify and prove a geometric or physical moduli-count correspondence with the same Hurwitz weights.

\section{Completion Estimate}
\noindent\textbf{COMPLETION: 0\%}
The calculation is an elementary benchmark or conditional reduction; it establishes no claimed new portion of the complete source target.

\subsection{Sources}
\begin{itemize}
\item[\hypertarget{source:20002511:1}{S1}] ULAM.AI and the UnsolvedMath Contributors, UnsolvedMath, record 20002511 (AIM-PHYSICS-0019). Catalogue curation under CC BY 4.0; definitions and English exposition added locally.
\url{https://huggingface.co/datasets/ulamai/UnsolvedMath}.
\item[\hypertarget{source:20002511:2}{S2}] American Institute of Mathematics, Gromov-Witten invariants and number theory, source item 4.2.
\url{http://aimpl.org/gromwitnumthry/4/}.
\item[\hypertarget{source:20002511:3}{S3}] Bringmann and Kane, work linked by the original source; not independently reviewed in this batch.
\url{https://arxiv.org/abs/1305.0112}.
\item[\hypertarget{source:20002511:4}{S4}] Bringmann and Murthy, work linked by the original source; not independently reviewed in this batch.
\url{https://arxiv.org/abs/1208.3476}.
\item[\hypertarget{source:20002511:5}{S5}] Bringmann and Manschot, work linked by the original source; not independently reviewed in this batch.
\url{https://arxiv.org/abs/1304.7208}.
\end{itemize}
\label{end:20002511}
\end{document}
